%%%%%%%%%%%%%%%%%%%%%%%%%%%%%%%%%%%%%%%%
%% MCM/ICM LaTeX Report
%% 2026 MCM Problem C
%%%%%%%%%%%%%%%%%%%%%%%%%%%%%%%%%%%%%%%%
\documentclass[12pt]{article}
\usepackage{geometry}
\geometry{left=1in,right=0.75in,top=1in,bottom=1in}

\newcommand{\Problem}{C}
\newcommand{\Team}{1111111}
\newcommand{\figdir}{../figures}
\newcommand{\tabledir}{../tables}

\usepackage{newtxtext}
\usepackage{amsmath,amssymb,amsthm}
\usepackage{newtxmath}
\usepackage[pdftex]{graphicx}
\usepackage{xcolor}
\usepackage{fancyhdr}
\usepackage{booktabs}
\usepackage{array}
\usepackage{float}
\usepackage{caption}
\usepackage{subcaption}
\usepackage{enumitem}
\usepackage[hidelinks]{hyperref}

\lhead{Team \Team}
\rhead{}
\cfoot{}
\setlength{\headheight}{14.5pt}
\setlength{\parskip}{0.35em}
\setlength{\parindent}{1.2em}
\setlist[itemize]{leftmargin=1.5em,itemsep=0.1em,topsep=0.2em}
\setlist[enumerate]{leftmargin=1.7em,itemsep=0.1em,topsep=0.2em}

\captionsetup{font=small,labelfont=bf}

\newtheorem{definition}{Definition}
\newtheorem{assumption}{Assumption}

\begin{document}
\pagenumbering{gobble}
\graphicspath{{../}{../figures/concept/}{../figures/generated/}}
\DeclareGraphicsExtensions{.pdf,.jpg,.jpeg,.png}

\thispagestyle{empty}
\vspace*{-16ex}
\centerline{\begin{tabular}{*3{c}}
  \parbox[t]{0.3\linewidth}{\begin{center}\textbf{Problem Chosen}\\ \Large \Problem\end{center}}
  & \parbox[t]{0.3\linewidth}{\begin{center}\textbf{2026\\ MCM/ICM\\ Summary Sheet}\end{center}}
  & \parbox[t]{0.3\linewidth}{\begin{center}\textbf{Team Control Number}\\ \Large \Team\end{center}} \\
  \hline
\end{tabular}}

\begin{center}
\Large \textbf{Data With the Stars: Inferring Hidden Fan Support Under Voting-Rule Constraints}
\end{center}

\noindent
Dancing with the Stars combines expert evaluation with audience voting, but the fan-vote component is secret. This makes Problem C fundamentally different from a direct prediction problem: the observed eliminations do not reveal a unique set of fan votes. Our solution therefore treats fan support as a latent variable. Instead of pretending to know exact vote totals, we reconstruct weekly competitions, translate observed eliminations into feasible fan-support constraints, estimate maximum-entropy fan-share distributions, and use those estimates to compare voting rules and propose a producer-facing rule.

\noindent
Using the official data set of 421 contestants across 34 seasons, we rebuilt 2,777 active contestant-week records and modeled 261 elimination weeks. The resulting latent-support model exactly reproduced the observed modeled eliminations by construction and confirmed that the underlying fan-vote uncertainty is large. This non-identifiability is the central practical finding: many fan-support profiles can explain the same observed elimination. It means that a good voting system should not merely chase an estimated ``true'' fan vote, but should be robust to close and ambiguous cases.

\noindent
Our counterfactual simulation shows that rank and percent rules disagree in 26.4\% of modeled elimination weeks, while estimated fan support changes the judge-only bottom set in 63.6\% of modeled weeks. We also examine controversial contestants, professional-dancer and celebrity-category effects, and the institutional role of bottom-two judge saves. These results support a hybrid recommendation: use a percent-based combined score for transparency, but introduce a gray-zone judge save when the bottom-two combined margin is small. In our simulation, the gray-zone trigger activates in 61.7\% of modeled elimination weeks, making close calls explicit rather than hiding them inside opaque vote totals.

\noindent
For producers, the recommendation is simple: keep fan influence visible, avoid overclaiming exact fan support, and reserve judge discretion for statistically fragile bottom-two cases. The proposed system is more explainable than a pure rank system and more stable than treating a tiny combined-score gap as decisive.

\clearpage
\pagestyle{fancy}
\pagenumbering{arabic}
\setcounter{page}{1}
\rhead{Page \thepage\ }

\section{Introduction}

\section{Data Audit and Preprocessing}

\section{Voting Rule Formalization}

\section{Fan Vote Estimation}

\section{Model Validation and Consistency}

\section{Rank vs. Percent Counterfactual Analysis}

\section{Controversial Contestants}

\section{Professional Dancer and Celebrity Effects}

\section{Proposed Voting System}

\section{Memo to DWTS Producers}

\section{Strengths, Limitations, and Extensions}

\begin{thebibliography}{9}

\bibitem{comapProblemC}
COMAP.
\newblock 2026 MCM Problem C: Data With the Stars.
\newblock Mathematical Contest in Modeling, 2026.

\bibitem{comapInstructions}
COMAP.
\newblock Contest Rules, Registration and Instructions for MCM/ICM.
\newblock \url{https://contest.comap.com/undergraduate/contests/mcm/instructions.php}, 2026.

\bibitem{shannon1948}
C. E. Shannon.
\newblock A mathematical theory of communication.
\newblock \emph{Bell System Technical Journal}, 27:379--423, 623--656, 1948.

\bibitem{boyd2004}
S. Boyd and L. Vandenberghe.
\newblock \emph{Convex Optimization}.
\newblock Cambridge University Press, 2004.

\bibitem{diamond2016}
S. Diamond and S. Boyd.
\newblock CVXPY: A Python-embedded modeling language for convex optimization.
\newblock \emph{Journal of Machine Learning Research}, 17(83):1--5, 2016.

\bibitem{scipy2020}
P. Virtanen et al.
\newblock SciPy 1.0: Fundamental algorithms for scientific computing in Python.
\newblock \emph{Nature Methods}, 17:261--272, 2020.

\bibitem{statsmodels2010}
S. Seabold and J. Perktold.
\newblock Statsmodels: Econometric and statistical modeling with Python.
\newblock Proceedings of the 9th Python in Science Conference, 2010.

\end{thebibliography}

\end{document}
